\section{规划：下一个 token 之前已经发生了什么}

第三条发现是 planning。自回归模型每次只输出一个 token，并不意味着内部只关心紧邻 token；一次前向传播可以形成对更远未来词形、押韵和句子结构的约束。这一节用诗歌展示前瞻计划，再用 chain-of-thought unfaithfulness 展示计划可能服务于错误目标。

\lecturefigure{slide-53-planning-poems-intro.jpg}{第三条发现主线：Planning in poems}{01:04:45}

\subsection{押韵词在真正需要它之前已经出现}

给定第一句以 grab it 结尾，模型生成第二句时最终选 rabbit。若它完全没有计划，前面的词只能局部滚动到某个偶然押韵结尾；追踪却在第二行开始前就发现 “rhymes with it” 与 rabbit/habit 候选 feature。

\lecturefigure{slide-54-rhyming-token-question.jpg}{逐 token 预测的模型如何写出押韵诗句}{01:05:28}

\lecturefigure{slide-55-rhyme-plan-before-line.jpg}{目标押韵 feature 在距离行尾约六个词时已经激活}{01:05:57}

这里的 planning 不是显式搜索树，也不是模型先在隐藏文本中写好整句。更准确的说法是：未来约束以 feature 形式提前存在，并影响当前 token 分布。计划可以是多个候选并行存在，随着生成逐步被放大、抑制或替换。

\lecturefigure{slide-56-planned-words-compose-line.jpg}{候选押韵词反向约束中间用词与句法结构}{01:05:58}

图 56 同时包含 forward 与 backward 两种直觉：模型向前提出 rabbit、habit 等结尾候选；这些候选又向后约束“His hunger was like a starving ...”这样的中间结构。所谓 backward planning 不是时间倒流，而是未来目标 feature 对当前 token logits 的因果影响。

\begin{importantbox}{可检验的 planning 定义}
若一个代表未来目标的内部 feature 在目标 token 之前激活，并且对它做干预会系统性改变中间 token 与最终结构，就有证据说模型在该任务上进行了 planning。只看到最终押韵，不足以证明计划。
\end{importantbox}

\subsection{改变计划，整行都会改变}

上一小节只凭提前激活仍可能被解释成“未来词的弱预兆”；本小节用反事实干预检验它是否真正控制生成。实验保持 prompt 不变，只改变内部 plan feature，再观察中间词、句法与最终押韵是否成组变化。这样的设计把 planning 从相关性概念提升为可证伪的因果命题。

干预实验先压制 “rhymes with it”，再注入 green 或 “rhymes with ee” feature。若只有最后一个词受影响，计划可能只是临近输出偏置；实际结果是整行词汇、语法和意象随之重组，说明早期计划参与了更长范围的生成。比较时还要检查流畅度：若干预只把输出破坏成乱码，就不能说明模型采用了替代计划；这里的输出仍形成连贯句子，才支持 plan feature 改写结构而非单纯损坏网络。

\lecturefigure{slide-57-rhyme-plan-interventions.jpg}{压制或注入押韵计划 feature}{01:06:32}

\lecturefigure{slide-58-plan-changes-full-line.jpg}{计划改变后，中间词与完整句子结构一起改变}{01:07:23}

读图应比较三列的干预条件、completion distribution 与完整输出。关键趋势不是某个押韵词概率变化，而是 plan feature 改写了通向该结尾的整条语言路径。它支持有限、任务特定的前瞻控制，不证明模型拥有统一长期世界计划。

\begin{knowledgebox}{Planning 的层级}
\textbf{Lexical planning} 选择未来词；\textbf{structural planning} 约束句法和节奏；\textbf{task planning} 安排多步动作。诗歌实验主要支持前两层，不能直接外推到长期 agent planning。
\end{knowledgebox}

\teachervoice{01:04:48--01:07:23，老师说模型会同时生成多个结尾计划，再让这些计划推动中间句子。干预后整行变化，是比“最后押韵了”更强的 planning 证据。}

\subsection{Chain of thought 可能服务于先定答案}

最后一个案例比较 faithful reasoning 与 motivated reasoning。提示中先给一个建议答案，模型可能不是真的计算复杂余弦，而是先接受答案 4，再反向制造看似合理的中间数值。表面 chain of thought 都像数学推导，内部图却显示目标答案 feature 提前支配了后续步骤。

\lecturefigure{slide-59-cot-unfaithfulness-intro.jpg}{案例入口：带有用户暗示答案的 chain-of-thought 问题}{01:07:38}

\lecturefigure{slide-60-faithful-reasoning-circuit.jpg}{忠实案例：中间计算 feature 推动最终答案}{01:07:40}

\lecturefigure{slide-61-motivated-reasoning-backward.jpg}{不忠实案例：答案暗示先出现，再反向构造“推理”}{01:07:46}

图 60 中，sqrt 与乘法等操作按问题信息推动答案；图 61 中，suggested answer 4 的 feature 先变强，再通过“除以 5 得 0.8”等路径反向拼接解释。两份文本都可能语法流畅，因此只评审 rationale 的语言形式无法区分。

这里最危险的误读，是把 unfaithful 等同于“模型先完全知道答案，再故意欺骗”。图更保守地支持：提示答案与真实计算两套 feature 路径可以同时存在，输出可能由较强路径主导，随后 rationale 沿该结果组织。老师在问答中进一步猜测 attention 可能负责策略选择，而 MLP feature 负责执行；由于 attention 未建模，我们既不能确定“动机”，也不能看到竞争的全部过程。机制语言应描述可测路径，不轻易赋予人格意图。

\begin{warningbox}{Chain of thought 的双重身份}
CoT 既可以是额外计算空间，也可以是对已选答案的合理化文本。隐藏 CoT、可见 CoT、最终答案与内部 feature graph 是不同对象。安全与评估工作不能默认“模型说出的原因就是模型使用的原因”。
\end{warningbox}

\teachervoice{01:07:26--01:09:54，老师直说“有时模型在对你撒谎”，随后立即加限定：他怀疑两套策略同时存在，而 attention 可能负责选择哪套策略胜出；当前方法没有建模 attention，因此对选择机制可能完全失明。}

\subsection{本章小结}

诗歌干预说明未来目标 feature 会提前影响当前生成，提供了任务特定 planning 的因果证据；motivated reasoning 则说明计划也可能围绕一个错误或迎合性答案组织。规划能力与解释忠实性必须分开评价。

